\documentclass[11pt]{article}
\usepackage[a4paper,margin=1in]{geometry}
\usepackage{amsmath,amssymb,mathtools,bm}
\usepackage{enumitem,booktabs,array}
\usepackage{tcolorbox}
\usepackage{hyperref}
\usepackage[protrusion=true,expansion=false]{microtype}
\hypersetup{colorlinks=true,linkcolor=blue,urlcolor=blue,citecolor=blue}
\setlist[itemize]{topsep=2pt,itemsep=2pt}
\setlist[enumerate]{topsep=2pt,itemsep=2pt}
\newtcolorbox{keybox}{colback=blue!3!white,colframe=blue!40!black,boxrule=0.4pt,arc=1mm}
\newtcolorbox{alertbox}{colback=red!3!white,colframe=red!40!black,boxrule=0.4pt,arc=1mm}
\newtcolorbox{answerbox}{colback=green!3!white,colframe=green!40!black,boxrule=0.4pt,arc=1mm}
\newtcolorbox{drillbox}{colback=yellow!5!white,colframe=orange!60!black,boxrule=0.4pt,arc=1mm}
\newcommand{\EV}{\mathbb{E}}
\newcommand{\Var}{\mathrm{Var}}
\newcommand{\Cov}{\mathrm{Cov}}
\newcommand{\blank}[1]{\underline{\hspace{#1}}}

\title{E11-02: Heider \& Ljungqvist (2015, JFE) \\
\large ``As Certain as Debt and Taxes: Estimating the Tax Sensitivity of Leverage from State Tax Changes'' --- Active Recall Workbook}
\author{Empirical Corporate Finance II --- Exam Prep}
\date{\today}
\begin{document}\maketitle

\begin{keybox}
\textbf{Paper in one sentence.} HL exploit plausibly-exogenous changes in US state corporate-income-tax rates as shocks to the tax benefit of debt, and show that firms raise (lower) their leverage in response to tax \emph{increases} (decreases), but the adjustment is highly asymmetric: firms react strongly to tax \emph{hikes} but barely at all to tax \emph{cuts}, producing a novel test of trade-off theory and evidence of frictions.

\textbf{Placement.} Standard modern identification paper in capital-structure research. Direct test of the tax-benefit channel in the trade-off theory.
\end{keybox}

\section*{Round 1: Complete Walk-Through}

\subsection*{1.1 Research question}
Trade-off theory predicts firms adjust leverage as tax benefits of debt change. Prior leverage-tax studies relied on marginal-tax-rate simulations (Graham 1996). Cleaner identification requires an exogenous shock. US state corporate tax changes provide many, largely-exogenous rate movements.

\subsection*{1.2 Setting and data}
\begin{itemize}
\item US state corporate income tax rate changes, 1989--2011.
\item 200+ tax events (hikes and cuts).
\item Compustat firms assigned to a ``home'' state (headquarters).
\item Treated firms: headquartered in states undergoing tax change.
\item Controls: firms in unaffected states (synthetic DiD or nearest-neighbor matched).
\end{itemize}

\subsection*{1.3 Empirical design}
Event-study DiD:
\[
\Delta\text{Leverage}_{it}=\alpha_i+\delta_t+\beta\cdot\Delta\text{Tax}_{s(i),t}+X_{it}'\gamma+\varepsilon_{it}.
\]
Separate estimates for hikes and cuts.

\subsection*{1.4 Headline results}
\begin{itemize}
\item Tax hikes: firms raise leverage by $\sim$0.4 pp per 1-pp tax increase over 2--3 years.
\item Tax cuts: leverage change is statistically zero.
\item Asymmetry robust across specifications.
\item Effect larger for financially-unconstrained firms (can easily access debt markets).
\end{itemize}

\subsection*{1.5 Mechanism}
\begin{itemize}
\item Tax benefit of debt rises with tax rate $\Rightarrow$ firms increase leverage post-hike.
\item Asymmetry reflects frictions: raising leverage is cheap if there is no distress premium; cutting leverage costly due to refinancing, signaling, debt-retirement costs.
\item Consistent with dynamic trade-off theory with adjustment costs.
\end{itemize}

\subsection*{1.6 Implications}
\begin{itemize}
\item Tax benefits matter for capital-structure decisions.
\item Leverage-tax elasticity is small but non-zero; previous small estimates do not necessarily refute trade-off theory.
\item Asymmetric adjustment highlights the importance of frictions in structural capital-structure models.
\end{itemize}

\subsection*{1.7 Caveats}
\begin{alertbox}
\begin{itemize}
\item State tax changes may correlate with local economic shocks.
\item Firms may have operations in multiple states; HQ assignment is imperfect.
\item Federal taxes dominate corporate-tax exposure; state-level is smaller.
\end{itemize}
\end{alertbox}

\section*{Round 2: Level 1 Fill-in}
\begin{enumerate}
\item Shock: US \blank{2cm} corporate tax changes.
\item Period: \blank{1cm}--\blank{1cm}.
\item Tax hike response: $+$\blank{1cm} pp leverage per 1-pp tax.
\item Tax cut response: \blank{2cm}.
\item Explanation: adjustment costs \& \blank{2cm} frictions.
\end{enumerate}

\section*{Round 3: Level 2 Prompts}
\begin{enumerate}
\item Describe the state-tax setting and identification.
\item Report main findings on hikes vs.\ cuts.
\item Explain the asymmetry through adjustment-cost theory.
\item Implications for trade-off theory.
\item Evaluate identification threats.
\end{enumerate}

\section*{Round 4: Minimal Scaffolding}
From a blank page: (i) question, (ii) shock, (iii) results, (iv) asymmetry, (v) implications.

\section*{Round 5: Blank Page}
Essay: taxes, adjustment costs, and capital structure.

\vspace{6pt}
\begin{drillbox}
\textbf{Speed Drill.} (1) Shock. (2) Sample years. (3) Hike effect. (4) Cut effect. (5) Interpretation.
\end{drillbox}

\section*{Quick Reference \& Answer Key}
\begin{answerbox}
\begin{itemize}
\item \textbf{Shock.} State corporate-tax changes.
\item \textbf{Hike.} Leverage $+0.4$ pp per 1-pp tax.
\item \textbf{Cut.} No response.
\item \textbf{Lesson.} Asymmetric, adjustment-cost-driven.
\end{itemize}
\end{answerbox}

\begin{keybox}
\textbf{30-second exam pitch.} Heider \& Ljungqvist (2015) use US state corporate-income-tax rate changes 1989--2011 --- 200+ largely-exogenous tax events --- to estimate the causal effect of tax rates on corporate leverage. In an event-study DiD comparing firms headquartered in treated vs.\ untreated states, they find that following a 1-pp tax hike, leverage rises $\sim$0.4 pp over 2--3 years; following a tax cut, the leverage response is statistically zero. This asymmetry is consistent with dynamic trade-off theory with adjustment costs: raising leverage is cheap while cutting leverage (refinancing, debt retirement, signaling) is costly. The paper provides one of the cleanest modern tests of the tax channel in capital structure, reconciles small-elasticity estimates with trade-off theory, and highlights the central role of financial frictions in structural capital-structure research.
\end{keybox}

\end{document}
